\documentclass[11pt,a4paper]{article}
\setlength{\oddsidemargin}{0in}
\setlength{\evensidemargin}{0in}
\setlength{\textwidth}{6.5in}
\setlength{\topmargin}{-0.5in}
\setlength{\textheight}{9.5in}
\usepackage{amsmath}
\usepackage{url}
\usepackage{booktabs}

\begin{document}

\noindent
\textbf{Farag Azzedin} (corresponding author) \\
Department of Information and Computer Science \\
King Fahd University of Petroleum and Minerals (KFUPM) \\
Dhahran, Saudi Arabia \\
Email: \texttt{fazzedin@kfupm.edu.sa} \\[1.5em]

\noindent
\textbf{To:} Editor-in-Chief \\
\textit{IEEE Transactions on Machine Learning in Communications and Networking (TMLCN)} \\[1.5em]

\noindent
\textbf{Subject:} Submission of Manuscript: ``\textit{Energy Allocation Between Computation and Wi-Fi Transmission in TinyML Federated Learning: A Measured Crossover}'' \\[1em]

\noindent
Dear Editor-in-Chief and Editorial Board,

\vspace{0.5em}
\noindent
We are pleased to submit our original research manuscript entitled \textbf{``Energy Allocation Between Computation and Wi-Fi Transmission in TinyML Federated Learning: A Measured Crossover''} for consideration for publication as a Regular Paper in \textit{IEEE Transactions on Machine Learning in Communications and Networking (TMLCN)}.

\vspace{0.5em}
\noindent
Federated learning (FL) algorithms deployed across wireless edge and Internet of Things (IoT) networks are overwhelmingly designed under the foundational assumption that wireless communication is the primary energy bottleneck, whereas on-device computation is negligible. While valid for smartphones connected via cellular uplinks, this premise fails in TinyML microcontroller networks where battery-powered nodes execute on-device gradient backpropagation over local RF transceivers.

\vspace{0.5em}
\noindent
To our knowledge, this manuscript presents the \textbf{first physically instrumented measurement of the computation-vs-communication energy crossover for federated learning over Wi-Fi}. Utilizing a high-side TI INA238 power monitor with sub-millisecond hardware GPIO phase markers on an Espressif ESP32-S3 microcontroller training Google Speech Commands keyword-spotting models ($3{,}120$-byte FP32 head updates), our work delivers the following contributions of strong interest to the TMLCN community:
\begin{enumerate}
    \item \textbf{Physical Crossover Characterization ($E^*$):} In a main campaign of 60 runs ($E \in \{1, 2, 5, 8, 10, 20\}$, ten repetitions each, one fixed protocol), we find that computation energy scales linearly as $\Delta E_{\text{comp}}(E) = 1.002 \cdot E - 0.010$\,mJ ($R^2 = 0.999$, intercept consistent with zero), while communication energy per client turn is about flat at $\Delta E_{\text{comm}} \approx 9.7$\,mJ. The marginal crossover is $E^* \approx 9.7$ epochs. A second session of 15 runs at the same setting gives $7.4$\,mJ and $E^* \approx 7.4$, so we report the crossover as a range of about $7.5$ to $10$ epochs.
    \item \textbf{Radio Tail Isolation:} A 45-run test with settling delays of $0$, $60$ and $200$\,ms shows that training that starts right after a radio burst absorbs about $0.5$\,mJ of tail energy ($p \approx 0.04$). A direct look at the traces shows a tail of $0.6$ to $0.8$\,mJ after each burst.
    \item \textbf{Analytical Modeling Benchmarks:} We evaluate five canonical theoretical baselines (Shannon capacity, dynamic CMOS switching $\kappa f^2 C$, combined analytical, datasheet currents, and communication dominance), exposing the physical failure mechanisms (MAC/CSMA protocol overhead, memory stalls, analog baseband quiescent floor, board regulation loss) that cause standard models to mispredict physical energy by $-99.4\%$ (Shannon), $-94.7\%$ (switching), $-89.0\%$ (combined) and about $+5000\%$ (communication dominance); the datasheet model lands within $10\%$ at its nominal point but has a wide range.
    \item \textbf{Duty-Cycled Sleep \& Networking Guidelines:} We derive a closed-form unified crossover model $E^*(P_s)$ incorporating device sleep power and channel contention. Because the board draws $561$\,mW merely remaining awake, duty-cycling into sleep between rounds shifts the crossover to $E^*_{\text{sleep}} \approx 11$ to $19$ epochs across empirical communication windows ($T_{\text{comm}} \in [85, 160]$\,ms).
    \item \textbf{Decoupling Crossover from Learning Optimality:} In a realistic non-IID FL benchmark ($K=20$ clients, Dirichlet skew $\alpha=0.3$, $10$\,dB noise), we show that the hardware crossover $E^*$ and the task learning optimum $E_{\text{opt}}$ are distinct quantities: $E=5$ achieves the lowest marginal active ($0.030$\,J) and gross awake energy ($0.218$\,J), whereas deep sleep pushes the optimum to $E=20$ ($0.52$\,J) due to round wake-up overheads.
    \item \textbf{Transparent Empirical Provenance:} We provide a complete, rigorous audit separating directly measured silicon metrics from model-assumed parameters and simulation traces.
\end{enumerate}

\vspace{0.5em}
\noindent
This manuscript is an original work, has not been published previously, and is not currently under consideration for publication elsewhere.

\vspace{0.5em}
\noindent
Thank you for your time and consideration of our work.

\vspace{1.5em}
\noindent
Sincerely,\\[1em]
\textbf{Farag Azzedin} (corresponding author) \\
Department of Information and Computer Science, KFUPM \\[1em]
\textbf{Md Abdus Shakur Siam} (first author) \\
Department of Electrical and Electronic Engineering, \\
Bangladesh University of Engineering and Technology (BUET)

\end{document}
